\section*{الفصل \ref{ch:b1:ocean-circulation} --- دوران المحيط}
% القيم: calc/ch29.py

\begin{solution}{exo:b1:ocean-circulation:1}
(أ)~عند \qty{4}{m/s} (الصيغة الأولى): $1000\,C_D = 0.29 + 0.775 + 0.481 = 1.546$،
$T = 1.3 \times \num{1.546e-3} \times 16 = \qty{0.0322}{Pa}$. عند \qty{7}{m/s}
(الصيغة الثانية): $1000\,C_D = 0.60 + 0.497 = 1.097$، $T = \qty{0.0699}{Pa}$. وعند
\qty{15}{m/s}: $1000\,C_D = 1.665$، $T = \qty{0.487}{Pa}$. (ب)~$0.487/0.0699 =
6.97$، مقابل $(15/7)^2 = 4.59$: معامل السحب نفسه يزداد مع الريح، لأن الريح الأقوى ترفع أمواجًا أخشن تمسك الهواء أفضل.
\end{solution}

\begin{solution}{exo:b1:ocean-circulation:2}
(أ)~نصف الكرة الشمالي، \ang{90} إلى يمين ريح تهبّ شرقًا: نحو الجنوب. (ب)~الريح تهبّ نحو الجنوب الغربي؛ ويمينها الشمال الغربي. (ج)~نصف الكرة الجنوبي، \ang{90} إلى يسار ريح تهبّ شرقًا: نحو الشمال. (د)~الريح تهبّ نحو الشمال؛ ويسارها الغرب، وهو عبر البحر على ساحل غربي. في الحالة (د) يشير النقل نحو البحر: فماء السطح يُستبدل من الأسفل، ويصعد الماء البارد على طول الساحل.
\end{solution}

\begin{solution}{exo:b1:ocean-circulation:3}
(أ)~$31.7 + 3.5 - 17.9 = \qty{17.3}{\sverdrup}$، مقابل \qty{17.1}{\sverdrup} المنشورة: تتّفق مجاميع الأجزاء الثلاثة مع الانقلاب المتوسط ضمن التقريب وعدم اليقين ($\pm\qty{3.2}{\sverdrup}$).
(ب)~فوق \qty{1100}{m}: $16.6 + 0.4 = \qty{17.0}{\sverdrup}$ نحو الشمال؛ وتحته: $12.0 + 5.9 = \qty{17.9}{\sverdrup}$ نحو الجنوب. يتعادل الاثنان تقريبًا، كما يجب حين لا يتراكم ماء شمال المقطع؛ والفرق الصغير ضمن عدم اليقين. (ج)~$\num{17.1e6} \times \num{3.156e7} = \qty{5.40e14}{m^3}$، أي \qty{5.40e5}{km^3} في السنة؛ و$\num{1.3324e9}/(\num{5.40e5}) \approx 2470$ سنة.
\end{solution}

\begin{solution}{exo:b1:ocean-circulation:4}
$D_E = 7.6 \times 10/\sqrt{\sin\varphi}$ و$V_0 = 0.0127 \times
10/\sqrt{\sin\varphi}$. عند \ang{15}: \qty{149}{m} و\qty{0.250}{m/s}
(\qty{2.5}{\percent} من الريح)؛ عند \ang{45}: \qty{90}{m} و\qty{0.151}{m/s}
(\qty{1.5}{\percent})؛ عند \ang{60}: \qty{82}{m} و\qty{0.136}{m/s}
(\qty{1.4}{\percent}). قرب خطّ الاستواء $\sin\varphi \to 0$ وتتباعد الصيغتان؛ والأهمّ أن $f \to 0$، فتنعدم قوة كوريوليس اللازمة لموازنة الاحتكاك، ويصبح الزمن $1/f$ الذي يحتاجه التوازن ليستقرّ أطول من أن تبقى الرياح ثابتة. وقصر إيكمان صيغته على $|\varphi| \geq
\ang{10}$.
\end{solution}

\begin{solution}{exo:b1:ocean-circulation:5}
(أ)~$f = 2\Omega\sin\ang{11} = \qty{2.783e-5}{s^{-1}}$؛ $M = 0.07/(\num{2.783e-5}) =
\qty{2515}{kg.m^{-1}.s^{-1}}$، \ang{90} إلى يمين ريح تهبّ غربًا: نحو الشمال. (ب)~$Q = (2515/1027) \times \num{5e6} = \qty{1.22e7}{m^3/s}$، أي \qty{12.2}{\sverdrup}. (ج)~يقع التقدير ضمن عدم يقين القياس المباشر ($12.0 \pm 5.5$) وقريبًا من القيمة المحسوبة من الرياح طويلة الأمد ($13.5 \pm 0.3$): فالنظرية تفسّر النقل المرصود.
\end{solution}

\begin{solution}{exo:b1:ocean-circulation:6}
(أ)~$f(\ang{2}) = \qty{5.09e-6}{s^{-1}}$؛ $M = 0.05/(\num{5.09e-6}) =
\qty{9824}{kg.m^{-1}.s^{-1}}$، أي \qty{9.57}{m^2/s}: عند \ang{2} شمالًا نحو الشمال (يمين ريح غربية)، وعند \ang{2} جنوبًا نحو الجنوب (يسارها). (ب)~عبر شريط طوله \qty{1}{m} على طول خطّ الطول، يغادر $2 \times 9.57 = \qty{19.1}{m^3/s}$ عبر خطّي العرض، المتباعدين \qty{444.8}{km}: $w = 19.1/(\num{4.448e5}) = \qty{4.30e-5}{m/s}$، أي نحو \qty{3.7}{m} في اليوم. (ج)~عند خطّ الاستواء $f = 0$ و$|\vec M| = T/|f|$ لا نهائي: لا يوجد توازن إيكمان المستقر هناك، والأرقام السابقة، الشديدة الحساسية لخطّ العرض المختار، ليست إلا رتبة قدر. ما يبقى هو الإشارة: فعلى كلّ جانب يتّجه النقل بعيدًا عن خطّ الاستواء، فيتباعد ماء السطح ويصعد الماء البارد على طوله.
\end{solution}

\begin{solution}{exo:b1:ocean-circulation:7}
(أ)~$f(\ang{15}) = \qty{3.775e-5}{s^{-1}}$: $M = 0.06/(\num{3.775e-5}) =
\qty{1590}{kg.m^{-1}.s^{-1}}$ نحو الشمال؛ $f(\ang{40}) =
\qty{9.375e-5}{s^{-1}}$: $M = \qty{1280}{kg.m^{-1}.s^{-1}}$ نحو الجنوب.
(ب)~بين خطّي العرض، المتباعدين \qty{2780}{km}، $\partial M_y/\partial y =
-(1590 + 1280)/(\num{2.78e6}) = \qty{-1.03e-3}{kg.m^{-2}.s^{-1}}$: وبحسب \cref{prop:b1:ocean-circulation:pumping}، $w_E = \qty{-1.005e-6}{m/s}$، أي سرعة هابطة قدرها \qty{31.7}{m} في السنة. (ج)~$\num{1.005e-6} \times \num{2.78e6} \times
\num{5e6} = \qty{1.40e7}{m^3/s}$، أي \qty{14.0}{\sverdrup} تُضخّ إلى الأسفل.
\end{solution}

\begin{solution}{exo:b1:ocean-circulation:8}
(أ)~$H = 1027 \times 3992 \times \num{17e6} \times 12 = \qty{8.36e14}{W}$، أي \qty{0.84}{PW}. (ب)~هي \qty{70}{\percent} من \qty{1.2}{PW} المقاسة: فالطبقة العليا الحقيقية أدفأ قرب السطح من متوسط \qty{15}{\celsius}، ويحمل \omterm{def:b1:ocean-circulation:gyre}{الدوّامة} حرارة أيضًا. (ج)~$1027 \times 3992 \times \num{17e6} \times 3 =
\qty{2.1e14}{W}$، أي \qty{0.21}{PW}، أي ربع (أ). فالماء الذي يصبح باردًا قبل عودته وحده يحمل حرارة كثيرة: ونقل حرارة الأطلسي يعتمد على تكوّن الماء العميق.
\end{solution}

\begin{solution}{exo:b1:ocean-circulation:9}
لتكن $x$ و$y$ و$1 - x - y$ كسور AAIW وNADW وAABW. توازنا الملح والحرارة، مقاسين من AABW:
\[
  -0.52\,x + 0.33\,y = 34.731 - 34.67 = 0.061,\qquad
  2.9\,x + 3.6\,y = 2.78 - 0.4 = 2.38 .
\]
محدِّد المصفوفة هو $-0.52 \times 3.6 - 0.33 \times 2.9 = -2.829$، ويعطي قاعدة كرامر $x = (0.061 \times 3.6 - 0.33 \times 2.38)/(-2.829) = 0.20$ و$y = (-0.52
\times 2.38 - 2.9 \times 0.061)/(-2.829) = 0.50$: فالعيّنة \qty{20}{\percent} AAIW، و\qty{50}{\percent} NADW و\qty{30}{\percent} AABW، أي نقطة داخل المثلث الذي تكوّنه المياه الثلاث.
\end{solution}

\begin{solution}{exo:b1:ocean-circulation:10}
(أ)~مع $W = u + iv = V_0e^{i\pi/4}e^{(1+i)az}$، وهو الذي جزؤه الحقيقي والتخيّلي هما صيغتا القضية، $W'' = (1+i)^2a^2W = 2ia^2W$. تجتمع المعادلتان في $A_z(u'' + iv'') = -fv + ifu = if(u + iv)$، أي $A_zW'' = ifW$، وهذا صحيح لأن $2ia^2A_z = if$ مع $a^2 = f/(2A_z)$. (ب)~$W'(0) =
(1+i)aV_0e^{i\pi/4} = \sqrt2\,aV_0\,e^{i\pi/2} = i\sqrt2\,aV_0$: إجهاد السطح $\rho A_zW'(0)$ يتّجه نحو الشمال، بمقدار $\rho A_z\sqrt2\,aV_0 = \rho A_z\sqrt2
\sqrt{f/(2A_z)}\;T/(\rho\sqrt{fA_z}) = T$. (ج)~$\int_{-\infty}^0 W\,\mathrm dz =
V_0e^{i\pi/4}/((1+i)a) = V_0/(\sqrt2\,a)$، عدد حقيقي: فليس للنقل جزء نحو الشمال، وجزؤه الشرقي هو $\rho V_0/(\sqrt2\,a) =
T/\sqrt{fA_z}\cdot\sqrt{A_z/f} = T/f$. وهذه هي النظرية: مقدار $T/f$، \ang{90} إلى يمين ريح تهبّ نحو الشمال.
\end{solution}

\begin{solution}{exo:b1:ocean-circulation:11}
(أ)~تغيّر خطوة نحو الشمال $\mathrm dy = a_\oplus\,\mathrm d\varphi$ قيمة $f = 2\Omega
\sin\varphi$ بمقدار $2\Omega\cos\varphi\,\mathrm d\varphi$، فيكون $\beta = 2\Omega
\cos\varphi/a_\oplus = 2 \times \num{7.292e-5} \times \cos\ang{30}/(\num{6.371e6}) =
\qty{1.98e-11}{m^{-1}.s^{-1}}$. (ب)~$R/\beta = \num{e-6}/(\num{1.98e-11}) =
\qty{5.04e4}{m}$، أي نحو \qty{50}{km}. (ج)~الداخل: $\num{30e6}/(\num{5e6} \times
1000) = \qty{6.0e-3}{m/s}$، أي \qty{6}{mm/s}؛ تيّار الحدّ: $\num{30e6}/(\num{5.04e4}
\times 1000) = \qty{0.60}{m/s}$، أسرع بمائة مرّة، وهي نسبة العرضين. ويبلغ تيّار الخليج، \qty{1.8}{m/s}، أسرع من ذلك: فالتيّار الحقيقي أضيق من هذا التقدير، وتجري تدفّقه قرب السطح، ويضيف دورانه المرتدّ إلى نقله.
\end{solution}

\begin{solution}{exo:b1:ocean-circulation:12}
(أ)~حيثما يكون البحر أعمق من \qty{1000}{m}، يحتوي العمود فوق \qty{1000}{m} على \qty{1}{km} من الماء على الأكثر، وفي غير ذلك أقلّ: فالحجم فوق \qty{1000}{m} لا يتجاوز $\num{361.84e6} \times 1 = \qty{3.618e8}{km^3}$، فيكون الحجم تحته لا يقل عن $\num{1.3324e9} - \num{3.618e8} = \qty{9.71e8}{km^3}$. (ب)~$\num{26e6} \times \num{3.156e7} = \qty{8.21e14}{m^3}$ في السنة،
\qty{8.21e5}{km^3}؛ و$\num{9.71e8}/(\num{8.21e5}) \approx 1200$ سنة، من الرتبة نفسها كألف سنة للناقل. (ج)~$(1/2)^{25/12.3} = 0.244$، أي نحو الربع؛
$\num{e-3} \times 25 \times \num{3.156e7} = \qty{7.9e5}{m}$، أي نحو \qty{790}{km}. وبعد 25 سنة انتشر التريتيوم لا أبعد مما كان سيحمله انجراف بطيء، بينما يبقى منه ربع يمكن كشفه.
\end{solution}

\begin{solution}{pb:b1:ocean-circulation:1}
\begin{enumerate}
  \item تهبّ ريح من الشمال نحو الجنوب، على طول الساحل.
  \item $1000\,C_D = 0.60 + 0.071 \times 8 = 1.168$: $C_D = \num{1.168e-3}$.
  \item $T = 1.3 \times \num{1.168e-3} \times 64 = \qty{0.0972}{Pa}$.
  \item $C_D = \num{1.452e-3}$ و$T = \qty{0.272}{Pa}$: يزداد الإجهاد بمقدار $2.80$، أكثر من $(12/8)^2 = 2.25$، لأن $C_D$ نفسه يزداد مع الريح.
  \item $\Omega = \qty{7.292e-5}{rad.s^{-1}}$؛ $f = 2\Omega\sin\ang{38} =
        \qty{8.979e-5}{s^{-1}}$.
  \item $2\pi/f = \qty{7.00e4}{s}$، أي \qty{19.4}{h}. تُدير قوة كوريوليس تيّارًا دورة كاملة في هذا الزمن؛ وقبل أن تهبّ الريح نحو هذا الزمن تقريبًا يتذبذب الماء بدلًا من أن يستقرّ في التوازن الثابت للنظرية.
  \item \ang{90} إلى يمين ريح تهبّ جنوبًا: نحو الغرب، أي نحو البحر المفتوح.
  \item $M = T/f = 0.0972/(\num{8.979e-5}) = \qty{1082}{kg.m^{-1}.s^{-1}}$.
  \item $1082/1027 = \qty{1.054}{m^2/s}$ لكل متر من الساحل.
  \item $D_E = 7.6 \times 8/\sqrt{\sin\ang{38}} = \qty{77.5}{m}$.
  \item $1.054/77.5 = \qty{0.0136}{m/s}$، أي \qty{1.36}{cm/s}.
  \item $V_0 = 0.0127 \times 8/\sqrt{\sin\ang{38}} = \qty{0.129}{m/s}$، بزاوية \ang{45} إلى يمين الريح: نحو الجنوب الغربي. وتعطي العوّامات نحو نصف قيمة إيكمان، أي نحو \qty{0.07}{m/s}.
  \item يخرج الماء من الساحل في \omterm{def:b1:ocean-circulation:ekman}{طبقة إيكمان} ولا يمكن تعويضه عبر الساحل، ولا من الأعلى: بل يعوّضه ماء يصعد من تحت الطبقة قرب الساحل، كما في المقطع العرضي المرسوم في \cref{sec:b1:ocean-circulation:upwelling}؛ وخطوط الحرارة المتساوية تنحني نحو الأعلى باتجاه الشاطئ.
  \item $1.054 \times 1000 \times 86400 = \qty{9.11e7}{m^3}$ في اليوم لكل كيلومتر من الساحل.
  \item $w = 1.054/(\num{2e4}) = \qty{5.27e-5}{m/s}$، أي \qty{4.55}{m} في اليوم.
  \item $77.5/4.55 = 17$ يومًا.
  \item يبرد الهواء من الأسفل بفعل الماء البارد: فتصبح الطبقة القريبة من السطح باردة ومستقرّة، وقد يتكثّف الرطوبة ضبابًا وسحبًا طبقية منخفضة، ومع غياب الحمل يكون المطر على الساحل قليلًا.
  \item الماء المصعَّد غني بالمغذّيات؛ وهي تغذّي العوالق النباتية، التي تغذّي العوالق الحيوانية والأسماك.
  \item ريح تهبّ نحو الشمال تدفع نقلًا \ang{90} إلى يمينها، نحو الشرق: يتراكم ماء السطح عند الساحل ويغوص (\omterm{def:b1:ocean-circulation:upwelling}{الهبوط})؛ فيتوقّف \omterm{def:b1:ocean-circulation:upwelling}{الصعود}، ويعود ماء سطح أدفأ إلى الساحل.
  \item $f(\ang{20}) = \qty{4.988e-5}{s^{-1}}$: $M = \qty{1948}{kg.m^{-1}.s^{-1}}$،
        أي $\sin\ang{38}/\sin\ang{20} = 1.80$ مرّة أكثر.
  \item في نصف الكرة الجنوبي يكون النقل \ang{90} إلى يسار الريح. ولكي يدفعه نحو الغرب، أي نحو البحر، يجب أن تهبّ الريح نحو الشمال، من الجنوب: وفي الحالتين تهبّ الريح على طول الساحل نحو خطّ الاستواء.
  \item $1000\,C_D = 0.29 + 0.62 + 0.308 = 1.218$؛ $T = \qty{0.0396}{Pa}$؛ $M =
        \qty{441}{kg.m^{-1}.s^{-1}}$؛ والنسبة $0.407$ في الحالتين.
  \item لكل متر من الساحل، $M c_p\,\Delta T = 1082 \times 3992 \times 6 =
        \qty{2.59e7}{W}$، أي \qty{25.9}{MW}؛ وعلى طول \qty{1000}{km}، \qty{2.59e13}{W}،
        أي \qty{25.9}{TW}.
  \item $\num{31.7e6}/1.054 = \qty{3.01e7}{m}$، أي نحو \qty{30000}{km}: يحرّك \omterm{def:b1:ocean-circulation:upwelling}{الصعود الساحلي} ماءً أقلّ بكثير مما يحرّكه تيّار حدّ غربي.
  \item $Q = 1.054 \times \num{e6} = \qty{1.05e6}{m^3/s}$: نقل \omterm{def:b1:ocean-circulation:upwelling}{الصعود} على طول \qty{1000}{km} من الساحل هو \qty{1.05}{\sverdrup}.
\end{enumerate}
\end{solution}
